\documentclass[11pt]{article}

\usepackage[margin=1in]{geometry}
\usepackage{booktabs}
\usepackage{listings}
\usepackage{xcolor}
\usepackage{hyperref}
\usepackage{amsmath}
\usepackage{parskip}
\usepackage{enumitem}

\definecolor{codebg}{rgb}{0.95,0.95,0.95}
\lstset{
    basicstyle=\ttfamily\small,
    backgroundcolor=\color{codebg},
    breaklines=true,
    frame=single,
    columns=fullflexible,
    language=C++,
    showstringspaces=false,
}

\newcommand{\field}[2]{\noindent\textbf{#1:} #2\\}

\title{Third-Party Notices}
\author{TOMOS}
\date{\today}

\sloppy
\emergencystretch=3em

\begin{document}
\maketitle

\noindent
This document enumerates the third-party software referenced or
incorporated in the development of TOMOS, for the purpose of
license-compliance review. For a general overview of the project, see the
project's \texttt{README.md}. The full, verbatim license text for every
entry below that is genuinely redistributed (compiled into a binary TOMOS
ships) is reproduced in the \texttt{licenses/} directory at the
repository root.

\section{Matrix Market I/O library (\texttt{mmio.c}/\texttt{mmio.h})}

\field{License}{Public domain}
\field{Author}{Roldan Pozo, NIST (\texttt{pozo@nist.gov}), Fri Aug 15 16:29:47 EDT 1997}
\field{Source}{\url{http://math.nist.gov/MatrixMarket}}

\noindent
\textbf{Description:} Vendored directly into this repository as
\texttt{src/io/mmio.c} and \texttt{src/io/mmio.h}, and compiled into
\texttt{tomos\_core}. This library provides the \texttt{.mtx} file reading
and writing functionality used throughout the project's benchmark and LSQR
tooling. Unlike every other entry in this document, it consists of vendored
source code residing within the repository, rather than a build-time fetch
or a link-only dependency. Both files contain the original NIST copyright,
permission, and warranty notice within their header comments.

\section{NVIDIA CUDA Toolkit (cuBLAS, cuSPARSE, cuSOLVER, cuBLASLt, cuRAND, CUDA Runtime)}

\field{License}{Proprietary; the NVIDIA CUDA Toolkit End User License Agreement (EULA). This is not an open-source license, and is neither BSD-3-Clause nor any OSI-approved license.}
\field{Source}{Not applicable; the toolkit is closed-source and distributed only as the CUDA Toolkit (\url{https://developer.nvidia.com/cuda-toolkit}).}

\noindent
\textbf{Description:} TOMOS's CUDA kernels (\texttt{tomos\_cuda}) and
solver code call directly into \texttt{cublas}, \texttt{cusparse},
\texttt{cusolver}, \texttt{cublasLt}, \texttt{curand}, \texttt{cudart}, and
\texttt{cudadevrt}. No component of the CUDA Toolkit is vendored or
redistributed; TOMOS links against it at build time only, referencing
whichever CUDA Toolkit installation is present on the end user's own
system (\texttt{CUDA\_HOME}). Anyone building or running TOMOS must
separately install the CUDA Toolkit and is thereby bound by NVIDIA's own
End User License Agreement.

\section{Thrust}

\field{License}{Apache License 2.0 (a limited number of utility classes are provided under the Boost Software License or a BSD license; see Thrust's own \texttt{LICENSE} file).}
\field{Copyright}{No single copyright line is present; copyright notices are given within the license terms themselves rather than as a standalone header.}
\field{Source}{\url{https://github.com/NVIDIA/thrust}, header-only, distributed as part of the CUDA Toolkit.}
\field{Full license text}{\texttt{licenses/thrust.txt}}

\noindent
\textbf{Description:} Thrust is header-only and bundled with the CUDA
Toolkit (located under \texttt{CUDA\_HOME/include}; it is neither fetched
nor vendored separately). It is included directly in
\texttt{src/kernels/cuda/sketch/countsketch\_cuda\_kernels.cu} and
\texttt{src/kernels/cuda/sketch/gaussian\_sketch\_cuda\_kernels.cu}, and is
compiled into \texttt{tomos\_cuda}. Unlike cuBLAS, cuSPARSE, and the other
components listed above, which are dynamically linked and never
incorporated into TOMOS's own binary, Thrust's template code is compiled
directly into \texttt{tomos\_cuda}'s object code. For this reason it is
listed as a separate entry rather than being subsumed under the
proprietary CUDA Toolkit license: it is a distinct, Apache-licensed
open-source component not covered by that End User License Agreement.
Thrust does not ship a \texttt{NOTICE} file, so no additional material
beyond the standard copyright and license text need be propagated.

\section{OpenBLAS}

\field{License}{BSD-3-Clause}
\field{Copyright}{Copyright (c) 2011-2014, The OpenBLAS Project}
\field{Source}{\url{https://github.com/OpenMathLib/OpenBLAS}, pinned at \texttt{v0.3.29}}
\field{Full license text}{\texttt{licenses/openblas.txt}}

\noindent
\textbf{Description:} OpenBLAS is fetched automatically at configure time
(\texttt{dependencies/CMakeLists.txt}, via CMake's \texttt{FetchContent}),
built from source, and linked into the shipped \texttt{tomos\_core} and
\texttt{tomos\_cpu} library binaries. It provides the CBLAS and LAPACKE
routines used by \texttt{src/core/precond/sketch\_qr.cpp} and
\texttt{src/kernels/cpu/dependencies/lapack.cpp}. This constitutes a
genuine redistribution, as OpenBLAS's compiled code becomes part of the
binaries that TOMOS ships.

\section{CUTLASS}

\field{License}{BSD-3-Clause}
\field{Copyright}{Copyright (c) 2017 - 2026, NVIDIA CORPORATION \& AFFILIATES}
\field{Source}{\url{https://github.com/NVIDIA/cutlass}}

\noindent
\textbf{Description:} CUTLASS is used in two distinct capacities. First, as
an external benchmark comparison: a small number of standalone files under
\texttt{benchmarks/kernels/cuda/matrix/} include CUTLASS headers directly
and instantiate \texttt{cutlass::gemm::device::Gemm} to obtain a
same-process throughput and correctness comparison. CUTLASS is not
vendored in this repository for this purpose; it is fetched separately by
whoever executes these scripts (via the \texttt{CUTLASS\_DIR} environment
variable), and the resulting comparison executables are built independently
of the \texttt{tomos\_core}, \texttt{tomos\_cpu}, and \texttt{tomos\_cuda}
library targets. Second, as a reference during kernel tuning: over the
course of a single GEMM kernel-tuning investigation confined to
\texttt{src/kernels/cuda/matrix/mm\_cuda\_kernels.cu}, specific technical
facts drawn from CUTLASS's source --- including a tile shape, a
lane-to-coordinate mapping formula, a PTX instruction choice, and a
pipeline-stage convention --- were used to inform the design of TOMOS's
own, independently written kernels. Each instance enumerated below is
already documented in the kernel file's own comments at the cited
location; the list that follows serves as an index into that existing
documentation rather than as new attribution text.

\subsection*{Design decisions informed by reading CUTLASS source}

\begin{enumerate}[leftmargin=*]

\item \textbf{Lane-to-(row,col) mapping.}
\textit{CUTLASS source cited:} \texttt{cutlass/layout/matrix.h},
\texttt{RowMajorInterleaved<2>::inverse()}.
\textit{Applied as:} a 2-line arithmetic formula replacing a plain
divide/modulo split.
\textit{Location:} \texttt{mm\_cuda\_kernels.cu:5378-5443}.

\item \textbf{ldmatrix PTX variant + offset toggle.}
\textit{CUTLASS source cited:}
\texttt{cutlass::gemm::warp::MmaTensorOpMultiplicandTileIterator},
\texttt{arch/memory\_sm75.h} --- via a standalone scratchpad that directly
\texttt{\#include}s real CUTLASS as ground truth.
\textit{Applied as:} swapped \texttt{.trans}$\to$non-transposed instruction,
XOR offset fix.
\textit{Location:} \texttt{mm\_cuda\_kernels.cu:15950-16229}.

\item \textbf{Instruction scheduling/interleaving.}
\textit{CUTLASS source cited:} \texttt{cuobjdump} SASS diff vs CUTLASS's
compiled mainloop; also noted CUTLASS has no \texttt{\_\_launch\_bounds\_\_}.
\textit{Applied as:} reordered load/compute source order; dropped a
register ceiling.
\textit{Location:} \texttt{mm\_cuda\_kernels.cu:1315-1346, 4242-4279,
6025-6041, 6669-6684, 25001-25059, 28281-28331}.

\item \textbf{Synchronous-load pipeline structure.}
\textit{CUTLASS source cited:} \texttt{mma\_pipelined.h}'s
\texttt{gemm\_iters}.
\textit{Applied as:} rewrote pipeline to match (no cp.async).
\textit{Location:} \texttt{mm\_cuda\_kernels.cu:1563-1586}.

\item \textbf{Multistage priming/stage-count convention.}
\textit{CUTLASS source cited:} \texttt{MmaMultistage}, confirmed
\texttt{Stages=3} via \texttt{cuobjdump} on the built reference binary.
\textit{Applied as:} wraparound-counter stage tracking; Stages=3 variants.
\textit{Location:} \texttt{mm\_cuda\_kernels.cu:2025-2040, 15646-15663,
20780-20818, 24542-24567, 30591-30695}.

\item \textbf{Precomputed-stride pointer addressing.}
\textit{CUTLASS source cited:} \texttt{PredicatedTileIterator};
corroborated via register-count comparison.
\textit{Applied as:} replaced per-tile IMAD address recompute, many
kernels.
\textit{Location:} \texttt{mm\_cuda\_kernels.cu:2242-2263, 2486-2514,
2742-2746, 5595-5606, 19434-19485, 46109-46120, 46765-46778}.

\item \textbf{Accumulator/warp-tile sizing.}
\textit{CUTLASS source cited:} SASS diff: CUTLASS's 128-float/thread,
square 64$\times$64 tile.
\textit{Applied as:} BN changed 160$\to$128 for an exact match.
\textit{Location:} \texttt{mm\_cuda\_kernels.cu:18105-18186}.

\item \textbf{A-load thread$\to$address mapping (fixed an L2 bug).}
\textit{CUTLASS source cited:} \texttt{PredicatedTileAccessIterator} +
\texttt{cp\_async} --- wired up via real CUTLASS headers in a standalone
repro.
\textit{Applied as:} per-thread K/M mapping rewritten; fixed a confirmed
1.5x L2-sector excess.
\textit{Location:} \texttt{mm\_cuda\_kernels.cu:18414-18609}.

\item \textbf{B-operand load mechanism.}
\textit{CUTLASS source cited:} ``zero STS, both operands via cp.async''.
\textit{Applied as:} converted B load to direct cp.async; $\sim$2x win.
\textit{Location:} \texttt{mm\_cuda\_kernels.cu:18610-18648}.

\item \textbf{A-layout swizzle formula.}
\textit{CUTLASS source cited:} \texttt{tensor\_op\_multiplicand\_sm75.h}'s
\texttt{TensorOpMultiplicand::operator()}.
\textit{Applied as:} corrected an earlier self-invented, wrong swizzle.
\textit{Location:} \texttt{mm\_cuda\_kernels.cu:23095-23167}.

\item \textbf{Epilogue vectorization width.}
\textit{CUTLASS source cited:} matching config's
\texttt{LinearCombination<..., 128/sizeof\_bits<ElementC>>}.
\textit{Applied as:} 8-half/float4 epilogue transfer.
\textit{Location:} \texttt{mm\_cuda\_kernels.cu:19885-20377}.

\item \textbf{Serpentine mma.sync accumulate order.}
\textit{CUTLASS source cited:} \texttt{mma\_tensor\_op.h}'s
\texttt{MmaTensorOp::operator()}.
\textit{Applied as:} reordered 32 mma.sync calls; \texttt{.reuse}
75\%$\to$97\%.
\textit{Location:} \texttt{mm\_cuda\_kernels.cu:27785-27858}.

\item \textbf{Full reference tile/warp/pipeline shape.}
\textit{CUTLASS source cited:} \texttt{cutlass\_fp16\_gemm\_ref.cu}'s
actual config (128$\times$256$\times$32 tile, 64$\times$64$\times$32 warp, 3
stages).
\textit{Applied as:} built from scratch together as one kernel.
\textit{Location:} \texttt{mm\_cuda\_kernels.cu:16838-16945}.

\item \textbf{B-swizzle (bigwarp family).}
\textit{CUTLASS source cited:} \texttt{tensor\_op\_multiplicand\_sm75.h},
\texttt{RowMajorTensorOpMultiplicandCongruous<16,64>}.
\textit{Applied as:} XOR partition swizzle for B.
\textit{Location:} \texttt{mm\_cuda\_kernels.cu:47478-47528}.

\item \textbf{Warp-grid arrangement.}
\textit{CUTLASS source cited:} CUTLASS's own
\texttt{examples/14\_ampere\_tf32\_tensorop\_gemm}.
\textit{Applied as:} 2$\times$2 warp grid, low-occupancy variant.
\textit{Location:} \texttt{mm\_cuda\_kernels.cu:44714-44748}.

\item \textbf{Register-level A/B double-buffering + spread cp.async
issuance.}
\textit{CUTLASS source cited:} \texttt{mma\_multistage.h} --- comment
explicitly says ``read verbatim via curl''.
\textit{Applied as:} ported into the ``bigwarp'' kernel family (one
sub-piece reverted after \texttt{compute-sanitizer} found a data race).
\textit{Location:} \texttt{mm\_cuda\_kernels.cu:44977-45404, 45530-45826,
46109-46778, 47100-47898}.

\item \textbf{Epilogue shared-memory staging.}
\textit{CUTLASS source cited:}
\texttt{epilogue/threadblock/epilogue.h}.
\textit{Applied as:} stage C through shared memory before coalesced global
write.
\textit{Location:} \texttt{mm\_cuda\_kernels.cu:15111-15136}.

\end{enumerate}

None of the instances enumerated above incorporate CUTLASS's source text
directly into TOMOS. In each case, a specific technical fact --- a formula,
a configuration parameter, or an instruction choice --- was identified in
CUTLASS's source and subsequently re-implemented in TOMOS's own,
independently structured kernel code, using TOMOS's own naming
conventions. Whether this practice constitutes the creation of a
derivative work under copyright law is a legal determination that falls
outside the scope of this document.

The same source file contains a small number of additional references to
``CUTLASS'' that invoke it as general design inspiration without citing a
specific header, class, or formula (for example, a comment describing a
``cutlass-style'' warp-tiling approach, or a reference to a ``real CUTLASS
technique'' in connection with an experiment that was never incorporated
into the shipped codebase). These references are omitted from the list
above because they do not cite a specific source; documentation of them is
available upon request.

\section{Google Test (gtest)}

\field{License}{BSD-3-Clause}
\field{Copyright}{Copyright 2008, Google Inc.}
\field{Source}{\url{https://github.com/google/googletest}, pinned at commit \texttt{03597a01ee50ed33e9dfd640b249b4be3799d395} (\texttt{test/CMakeLists.txt})}
\field{Full license text}{\texttt{licenses/googletest.txt}}

\noindent
\textbf{Description:} Google Test is fetched automatically at configure
time via CMake's \texttt{FetchContent}, built from source, and linked into
every test executable under \texttt{test/} (\texttt{GTest::gtest\_main}).
This is the same redistribution situation as OpenBLAS: gtest's compiled
code becomes part of the shipped test binaries.

\section{Google Benchmark}

\field{License}{Apache License 2.0}
\field{Copyright}{No single copyright line is present; the project's \texttt{AUTHORS} file attributes the work to multiple individual and organizational contributors, originally including Google Inc.}
\field{Source}{\url{https://github.com/google/benchmark}, pinned at \texttt{v1.9.1} (\texttt{benchmarks/CMakeLists.txt})}
\field{Full license text}{\texttt{licenses/google-benchmark.txt}}

\noindent
\textbf{Description:} Google Benchmark is fetched automatically at
configure time via CMake's \texttt{FetchContent}, built from source, and
linked into every \texttt{bench\_*} executable under \texttt{benchmarks/}
(\texttt{benchmark::benchmark}). Apache License 2.0 differs from every
other license family represented in this document in that it includes an
explicit patent grant and a NOTICE-file-propagation clause (\S4(d));
however, the google/benchmark project does not ship a \texttt{NOTICE}
file, so no additional material beyond the standard copyright and license
text need be propagated.

\end{document}
